\documentclass[10pt,a4paper]{amsart}
\usepackage[margin=19mm]{geometry}
\usepackage{amsmath,amssymb,mathtools}
\usepackage[scr=rsfs,cal=euler]{mathalfa}
\usepackage{xfrac}
\usepackage[unicode,hidelinks,bookmarks=false]{hyperref}
\newcommand{\ie}{i.e.\ }
\newcommand{\cf}{cf.\ }
\newcommand{\resp}{respectively\ }
\newcommand{\Subsection}{\subsection}
\providecommand{\nobreakdash}{\nobreak\hbox{-}}
\setlength{\emergencystretch}{2em}
\multlinegap0pt
\usepackage{amsthm}
\usepackage{mathtools}
\usepackage{amsmath,amssymb}
\usepackage{enumitem}
\usepackage{dsfont}
\usepackage{bbm}
\usepackage{mathrsfs}
\usepackage{xcolor}
\usepackage{graphicx}
\usepackage{cleveref}
\def\AA{{\mathcal{A}}}
\providecommand{\Exp}[1]{{ \mathbb{E}}\!\left[#1\right]}
\providecommand{\Expu}[2]{{\mathbb{E}}_{#1}\!\left[#2\right]}
\def\Fopt{{F^*}}
\def\HH{{\mathcal{H}}}
\def\RR{{\mathbb{R}}}
\def\Reg{{\mathcal{R}}}
\def\RegMax{{\mathcal{R}_{\mathsf{UB}}}}
\def\SS{{\mathcal{S}}}
\def\XX{{\Theta}}
\def\aa{{a}}
\DeclarePairedDelimiter{\abs}{\lvert}{\rvert}
\providecommand{\card}[1]{\left|#1\right|}
\def\dOpt{{d^*}}
\providecommand{\define}{\coloneqq}
\def\gradPsi{{\ell_{\psi}}}
\providecommand{\innerEmpty}{\,\cdot\,}
\def\jacPsi{{L_{\psi}}}
\DeclarePairedDelimiter{\norm}{\lVert}{\rVert}
\def\nuOpt{{\nu^{*}}}
\providecommand{\pen}[1]{{\plus{#1}^2}}
\providecommand{\penGenNot}{{\mathcal{L}}}
\DeclarePairedDelimiter{\plus}{[}{]_{+}}
\def\ss{{s}}
\def\xx{{\theta}}
\def\xxOpt{\xx^{*}}
\theoremstyle{plain}
\newtheorem{theorem}{Theorem}[section]
\newtheorem{proposition}[theorem]{Proposition}
\newtheorem{lemma}[theorem]{Lemma}
\newtheorem{corollary}[theorem]{Corollary}
\theoremstyle{definition}
\newtheorem{definition}[theorem]{Definition}
\newtheorem{example}[theorem]{Example}
\newtheorem{assumption}[theorem]{Assumption}
\theoremstyle{remark}
\newtheorem{remark}[theorem]{Remark}
\newtheorem{observation}[theorem]{Observation}
\title[Global optimality for constrained exploration]{Global Optimality for Constrained Exploration via Penalty Regularization: the smoothed-entropy regularity and dual translation bundle (Lemma C.3)}
\author[F. Wolf, I. Fatkhullin, N. He]{Florian
Wolf, Ilyas Fatkhullin, Niao He}
\date{Source: arXiv:2604.28144v2 (CC BY 4.0)}
\begin{document}
\maketitle
\vspace{1em}
\noindent\textit{Source and licence.} Global Optimality for Constrained Exploration via Penalty Regularization. Version arXiv:2604.28144v2 (CC BY 4.0), distributed under the Creative Commons Attribution 4.0 International licence, http://creativecommons.org/licenses/by/4.0/, as recorded in the arXiv OAI licence metadata for this exact version and shown on the abstract page https://arxiv.org/abs/2604.28144v2. This item is an editorial re-presentation of the paper's smoothed-entropy regularity and dual-translation lemma together with the paper's own proof of it and the source-local material that proof uses. A full rights notice is delivered at licenses/arxiv-2604.28144-CC-BY-4.0.txt.
\noindent\textit{Editorial scope.} Counted unit: the paper's smoothed-entropy regularity and dual-translation lemma, printed as Lemma C.3 (source0.tex lines 2454--2534), delivered together with the paper's complete original proof of it (source0.tex lines 2535--2586) as the single primary proof body. No mathematics is written, repaired or re-derived: statement and proof are the paper's own text line for line, in the paper's own notation ($\HH$, $\HH_\sigma$, $\lambda$, $\sigma$, $\nuOpt$, $\Reg$, $r$, $a_\sigma$). Required context retained uncounted, because the counted proof is the paper's own assembly over it: the constrained maximum-entropy problem itself (source0.tex lines 856--873); the three structural assumptions it invokes (lines 1133--1141, 1181--1192 and 1255--1267); and the smoothness lemma whose constants the proof substitutes (lines 2026--2075). Every cross-reference inside the retained spans resolves inside the delivered document, and no citation command occurs in any retained span, so the delivered document carries no bibliography and no reference or citation command of the delivered text was rewritten or adapted. Printed numbers are the paper's own: the wrapper restores the paper's section counter, its single shared section-relative theorem counter and its global equation counter before the retained material, so the retained assumptions print as Assumptions 5.1, 5.3 and 5.5, the retained smoothness lemma as Lemma B.3, the counted statement as Lemma C.3, and the numbered displays of the retained spans carry the paper's own equation numbers (1) and (7) to (16). Omitted from this package and not redistributed: the paper's preamble and front matter as such (of which only the title and the author names are reproduced in the wrapper); the introduction and related work; the main convergence theorem and its appendix proof; the remaining lemmas, corollaries and algorithms; the experiments; and the closing sections. The paper's own bibliography and its bib data file are likewise not redistributed; since no retained span cites anything, the delivered document carries no bibliography. Scope note, disclosed: the counted lemma is the paper's regularity bundle that turns the abstract penalty convergence theorem into the concrete rates for the smoothed-entropy penalty, and the paper gives its proof in full in the retained range; the main convergence theorem is proved in the paper's appendix over a chain that reaches the whole appendix, so it is not the counted unit. Printed slips kept literal, among them the paper's own wording and its own choice of symbols. Layout and class: the paper's own class is the standard LaTeX article class with the publisher's conference style file for the submission, which is neither shipped with this package nor used to build it; the delivered wrapper is the lane's pinned amsart editorial wrapper at 10pt on A4, to which this package adds the paper's own theorem-like declarations with their shared section-relative counter and the shorthand macros its retained text uses, all copied verbatim from the paper's source and its own local style file. No publisher class file, no publisher style file and no upstream script is redistributed. Assets: no retained span of this package contains a figure environment or a graphics-inclusion command, no image or artwork file exists in or is redistributed with this package, and the manifest and the delivery evidence both carry an empty asset-dependency list. A full rights notice is delivered at licenses/arxiv-2604.28144-CC-BY-4.0.txt.
\setcounter{section}{1}
Intuitively, since the OM captures the state-action visitation frequencies
induced by a policy, it provides a principled foundation for defining
exploration as maximizing the entropy of the induced state-action
distribution.

\section{Problem Formulation}
\label{sec:ProblemFormulation}
For a convex, $L$-smooth constraint $\Reg: \Lambda \to \RR$ formulated in
the state-action-occupancy measure, we formulate the following
constrained-max-entropy problem:
\begin{align}
 \max_{\theta \in \Theta} F_1(\theta)  \define \HH(\lambda^{\pi_\theta})
 = - \Expu{(\ss, \aa) \sim \lambda^{\pi_\theta}}{\log \lambda^{\pi_\theta}(\ss, \aa)} \;\;
 \mathrm{s.t.} \;\; F_2(\theta)        \define
 \Reg(\lambda^{\pi_\theta}) \leq 0.
 \tag{CME}
 \label{eq:MainProblemRL}
\end{align}

\setcounter{section}{5}
\setcounter{theorem}{0}
\begin{assumption}[Softmax-Parametrization]\label{ass:Softmax}
 The function $\psi : \SS \times \AA \times \Theta \to \RR$
 satisfies $\psi \in C^2(\SS, \AA, \Theta; \RR)$, and there exist
 constants $\gradPsi, \jacPsi > 0$ such that
 (i) $\max_{\ss\in\SS, \aa\in\AA} \sup_{\theta \in \Theta}
  \norm{\nabla_\theta \psi(\ss, \aa, \theta)} \leq \gradPsi$
 and (ii) $\max_{\ss\in\SS, \aa\in\AA} \sup_{\theta \in \Theta}
  \norm{\nabla^2_\theta \psi(\ss, \aa, \theta)} \leq \jacPsi$.
 \end{assumption}

\setcounter{theorem}{2}
\begin{assumption}[L-smoothness]
 \label{ass:Lsmoothness}
 Assume that for $\HH$ and $\Reg$ there exist
 $\ell_\lambda, L_\lambda, L_{\lambda, \infty} > 0$
 such that $G \in \{-\HH, \Reg\}$ satisfies for
 all $\lambda, \lambda^\prime \in \Lambda$:
 (i) $\norm{\nabla_\lambda G(\lambda)}_{\infty} \leq \ell_\lambda$,
 (ii) $\norm{\nabla_{\lambda} G(\lambda) - \nabla_\lambda G(\lambda^\prime)}_{\infty}
  \leq L_\lambda \norm{\lambda - \lambda^\prime}_2$, and
 (iii) $\norm{\nabla_{\lambda} G(\lambda) - \nabla_\lambda G(\lambda^\prime)}_{\infty}
  \leq L_{\lambda, \infty} \norm{\lambda - \lambda^\prime}_1$.
\end{assumption}

\setcounter{theorem}{4}
\begin{assumption}[Strong Duality]
 \label{ass:StrongDuality}
 We assume \emph{strong duality} holds for \eqref{eq:MainProblemRL},
 i.e. for $\xxOpt$ there exists a Lagrangian multiplier
 $\nuOpt \geq 0$ such that $\dOpt = d(\nuOpt) = -\Fopt$, where
 the \emph{(Lagrangian) dual function} is defined as
 the infimum of the \emph{Lagrangian} with respect to the primal variable $\xx$, that is,
 $d(\nu)\define\inf_{\xx\in\XX}L(\xx,\nu)
  =\inf_{\xx\in\XX}\bigl(-F_1(\xx)+\nu F_2(\xx)\bigr)$,
 and the corresponding \emph{dual problem} as
 $\dOpt \define \sup_{\nu \geq 0} d(\nu)
  = \sup_{\nu \geq 0} \inf_{\xx \in \XX} L(\xx, \nu)$.
\end{assumption}

\appendix
\setcounter{section}{2}
\setcounter{theorem}{2}
\begin{lemma}[Smoothness and Bounded Gradients of Penalty
  function in $\lambda$]
 \label{lem:AppendixSmoothnessBoundedGradientsPenaltyLambda}
 With $\penGenNot = \pen{\innerEmpty{}}$, under
 \Cref{ass:Softmax} and
 \Cref{ass:Lsmoothness}, the
 function $\lambda \mapsto P(\lambda)
  =-\HH(\lambda)+\beta\pen{\Reg(\lambda)}$
 fulfills the inequalities of \Cref{ass:Lsmoothness} with
 \begin{align*}
  \ell_{P,\lambda}
   & \define \ell_{\HH,\lambda}
  +2\beta\RegMax\ell_{\Reg,\lambda},                          \\
  L_{P,\lambda}
   & \define L_{\HH,\lambda}
  +2\beta\left(\RegMax L_{\Reg,\lambda}
  +\sqrt{\card{\SS}\card{\AA}}\,\ell_{\Reg,\lambda}^2\right), \\
  L_{P,\lambda,\infty}
   & \define L_{\HH,\lambda,\infty}
  +2\beta\left(\RegMax L_{\Reg,\lambda,\infty}
  +\ell_{\Reg,\lambda}^2\right),
 \end{align*}
 with the upper bound $\RegMax \define \max_{\lambda \in \Lambda} \abs{\Reg(\lambda)}$.
 In particular, the smoothness constants of the entropy function
 enter additively and are not multiplied by $\beta$.
\end{lemma}
\begin{proof}
 Writing $r(\lambda)=\plus{\Reg(\lambda)}$, the chain rule gives
 \begin{align*}
  \nabla_\lambda P(\lambda)
  =-\nabla_\lambda\HH(\lambda)
  +2\beta r(\lambda)\nabla_\lambda\Reg(\lambda).
 \end{align*}
 The gradient bound follows from $0\leq r\leq\RegMax$.
 The penalty-gradient difference is bounded by
 \begin{align*}
  2\beta r(\lambda)
  \norm{\nabla\Reg(\lambda)-\nabla\Reg(\lambda^\prime)}_\infty
  +2\beta\abs{r(\lambda)-r(\lambda^\prime)}
  \norm{\nabla\Reg(\lambda^\prime)}_\infty.
 \end{align*}
 The positive-part map is non-expansive, and hence
 $\abs{r(\lambda)-r(\lambda^\prime)}
  \leq\ell_{\Reg,\lambda}\norm{\lambda-\lambda^\prime}_1$.
 For the Euclidean bound, the assumed infinity-norm gradient bound gives
 $\norm{\nabla\Reg}_2\leq
  \sqrt{\card{\SS}\card{\AA}}\,\ell_{\Reg,\lambda}$.
 Combining these inequalities with the function-specific smoothness
 constants proves all three claims.
\end{proof}

\setcounter{section}{3}
\setcounter{theorem}{2}
\setcounter{equation}{6}
\begin{lemma}[Smoothed entropy, penalty constants, and dual translation]
 \label{lem:SmoothedEntropyProperties}
 Let $d\define\card{\SS}\card{\AA}$ and $\sigma\in(0,1]$.
 For every $\lambda\in\RR_+^d$ with $\norm{\lambda}_1\leq1$, the surrogate
 the smoothed surrogate $\HH_\sigma(\lambda)
  \define-\sum_{\ss,\aa}\lambda(\ss,\aa)
  \log\bigl(\lambda(\ss,\aa)+\sigma\bigr)$
 is concave and satisfies
 \begin{align}
  \nabla_i(-\HH_\sigma)(\lambda)
   & =\log(\lambda_i+\sigma)
  +\frac{\lambda_i}{\lambda_i+\sigma},
  \label{eq:SmoothedEntropyGradient}                  \\
  \nabla_{ii}^2(-\HH_\sigma)(\lambda)
   & =\frac{\lambda_i+2\sigma}{(\lambda_i+\sigma)^2}.
  \label{eq:SmoothedEntropyHessian}
 \end{align}
 Consequently, $-\HH_\sigma$ satisfies the three regularity
 inequalities used in \Cref{ass:Lsmoothness} with
 \begin{align}
  \ell_{\HH_\sigma,\lambda}
   & \define 1+\log\frac{2}{\sigma},
   & L_{\HH_\sigma,\lambda}
   & =L_{\HH_\sigma,\lambda,\infty}\define\frac{2}{\sigma}.
  \label{eq:SmoothedEntropyConstants}
 \end{align}
 In particular, \eqref{eq:SmoothedEntropyGradient} equals
 $\log\sigma$ at an empirical zero. Moreover, with
 \begin{align}
  a_\sigma\define\log(1+d\sigma)\leq d\sigma,
 \end{align}
 we have the uniform approximation bound
 \begin{align*}
  0\leq\HH(\lambda)-\HH_\sigma(\lambda)\leq a_\sigma.
 \end{align*}
 For $P_\sigma=-\HH_\sigma+\beta\pen{\Reg}$, its constants can
 therefore be chosen as
 \begin{align}
  \ell_{P_\sigma,\lambda}
   & \leq1+\log\frac{2}{\sigma}
  +2\beta\RegMax\ell_{\Reg,\lambda},
  \label{eq:SmoothedPenaltyGradientConstant} \\
  L_{P_\sigma,\lambda}
   & \leq\frac{2}{\sigma}
  +2\beta\left(\RegMax L_{\Reg,\lambda}
  +\sqrt d\,\ell_{\Reg,\lambda}^2\right),
  \label{eq:SmoothedPenaltyL2Constant}       \\
  L_{P_\sigma,\lambda,\infty}
   & \leq\frac{2}{\sigma}
  +2\beta\left(\RegMax L_{\Reg,\lambda,\infty}
  +\ell_{\Reg,\lambda}^2\right).
  \label{eq:SmoothedPenaltyL1Constant}
 \end{align}
 Finally, let $\lambda^*$ be an optimizer of the original (unsmoothed)
 problem \eqref{eq:MainProblemRL}, let $\nuOpt$ be the multiplier from
 \Cref{ass:StrongDuality}, and set $r(\lambda)=\plus{\Reg(\lambda)}$.
 Then the same original strong-duality assumption implies, that for every
 occupancy $\lambda=\lambda^{\pi_\theta}$ generated by some
 $\theta\in\Theta$, it holds
 \begin{align}
  -\HH_\sigma(\lambda)+\HH_\sigma(\lambda^*)
  +\beta r(\lambda)^2
  \geq\beta r(\lambda)^2-\nuOpt r(\lambda)-a_\sigma.
  \label{eq:SmoothedPenaltyDualLowerBound}
 \end{align}
 Thus, if $\lambda$ satisfies
 $\Exp{-\HH_\sigma(\lambda)+\HH_\sigma(\lambda^*)
   +\beta r(\lambda)^2}\leq\delta$ and
 $s\define\Exp{r(\lambda)}$, then
 \begin{align}
  s
   & \leq\frac{\nuOpt+
   \sqrt{\nuOpt^2+4\beta(\delta+a_\sigma)}}{2\beta},
  \label{eq:SmoothedConstraintTranslation}       \\
  \HH(\lambda^*)-\Exp{\HH(\lambda)}
   & \leq\delta+a_\sigma, \quad \text{and} \quad
  \Exp{\HH(\lambda)}-\HH(\lambda^*)
  \leq\nuOpt s.
  \label{eq:SmoothedObjectiveTranslation}
 \end{align}
\end{lemma}

\begin{proof}
 Differentiating each coordinate of
 $-\HH_\sigma(\lambda)=\sum_i\lambda_i\log(\lambda_i+\sigma)$ gives
 \eqref{eq:SmoothedEntropyGradient}--\eqref{eq:SmoothedEntropyHessian}.
 The Hessian is diagonal and positive semidefinite, and
 \begin{align*}
  0\leq\frac{x+2\sigma}{(x+\sigma)^2}\leq\frac{2}{\sigma}
  \qquad(x\geq0).
 \end{align*}
 The mean-value theorem then gives both smoothness inequalities.
 Since $x+\sigma\in[\sigma,2]$ for $x\in[0,1]$, the gradient formula gives
 $\norm{\nabla\HH_\sigma}_\infty\leq1+\log(2/\sigma)$.

 For the approximation bound, define $f_\sigma(x)=x\log(1+\sigma/x)$ for
 $x>0$ and $f_\sigma(0)=0$. Direct differentiation gives
 $f_\sigma^{\prime\prime}(x)=-\sigma^2/[x(x+\sigma)^2]\leq0$. For
 $g(m)=m\log(1+d\sigma/m)$, $g^{\prime}(m)=\log(1+t)-t/(1+t)\geq0$ with
 $t=d\sigma/m$. Thus, if $m=\norm{\lambda}_1>0$, concavity of $f_\sigma$ and
 $m\leq1$ give
 \begin{align*}
  \HH(\lambda)-\HH_\sigma(\lambda)
  =\sum_i f_\sigma(\lambda_i)
  \leq m\log\left(1+\frac{d\sigma}{m}\right)
  \leq\log(1+d\sigma)\leq d\sigma.
 \end{align*}
 The case $m=0$ follows by continuity. Substitution of
 \eqref{eq:SmoothedEntropyConstants} into
 \Cref{lem:AppendixSmoothnessBoundedGradientsPenaltyLambda} proves
 \mbox{\eqref{eq:SmoothedPenaltyGradientConstant}-
  \eqref{eq:SmoothedPenaltyL1Constant}}.

 Strong duality for the original problem gives, for every occupancy
 $\lambda=\lambda^{\pi_\theta}$ with $\theta\in\Theta$, the inequality
 \begin{align*}
  -\HH(\lambda)+\nuOpt\Reg(\lambda)
   & \geq d(\nuOpt)=-\HH(\lambda^*).
 \end{align*}
 Since $\Reg(\lambda)\leq r(\lambda)$, this implies the pointwise
 bound $\HH(\lambda)-\HH(\lambda^*)\leq\nuOpt r(\lambda)$.
 Writing $q_\sigma=\HH-\HH_\sigma\in[0,a_\sigma]$ therefore yields
 \begin{align*}
  \HH_\sigma(\lambda)-\HH_\sigma(\lambda^*)
  \leq\nuOpt r(\lambda)+a_\sigma,
 \end{align*}
 which is \eqref{eq:SmoothedPenaltyDualLowerBound}. Taking
 expectations and applying Jensen's inequality yields
 $\beta s^2-\nuOpt s\leq\delta+a_\sigma$; solving this quadratic proves
 \eqref{eq:SmoothedConstraintTranslation}. Finally, using
 $q_\sigma\in[0,a_\sigma]$ and the nonnegativity of the penalty gives the
 first inequality in \eqref{eq:SmoothedObjectiveTranslation}; the second
 is the original dual bound, concluding the proof.
\end{proof}

\end{document}
